% CSU Strategic Direction: Past Five Years to CSU Forward
% Detailed comparison report for FYE Skill Development & Career Readiness project

\documentclass[11pt,letterpaper]{article}

\usepackage[margin=1in]{geometry}
\usepackage{helvet}
\renewcommand{\familydefault}{\sfdefault}
\usepackage[T1]{fontenc}
\usepackage[utf8]{inputenc}
\usepackage{array}
\usepackage{tabularx}
\usepackage{booktabs}
\usepackage{url}

\pagestyle{plain}

\setlength{\parindent}{0pt}
\setlength{\parskip}{0.45em}

\makeatletter
\renewcommand\section{\@startsection{section}{1}{\z@}%
  {1.1em}{0.45em}{\large\bfseries}}
\renewcommand\subsection{\@startsection{subsection}{2}{\z@}%
  {0.85em}{0.3em}{\normalsize\bfseries}}
\makeatother

\newcommand{\sectionrule}{\par\vspace{0.2em}\noindent\rule{\textwidth}{0.5pt}\vspace{0.5em}}

\begin{document}

\begin{center}
  {\LARGE\bfseries CSU Strategic Direction Over the Past Five Years}\\[0.35em]
  {\large From Graduation Initiative 2025 to CSU Forward}\\[0.6em]
  {\normalsize A Detailed Comparison Report}\\[0.8em]
  {\small Prepared for: FYE Research \& Engagement Project}\\
  {\small Skill Development \& Career Readiness for FYE Student Employees}\\[0.35em]
  {\small Student: Pushpak Sunil Rane \quad|\quad Supervisor: Dani Fernandez}\\
  {\small California State University, Chico \quad|\quad Fall 2026}
\end{center}

\sectionrule

\begin{abstract}
\noindent
This report reviews how California State University (CSU) strategy evolved over roughly the last five years, culminating in the 2025 launch of \textit{CSU Forward: Thriving Students, Thriving University, Thriving California}. During most of this period, CSU's primary student-success agenda was Graduation Initiative 2025 (GI~2025), which focused on raising graduation rates and closing equity gaps. GI~2025 produced major gains---especially in four-year completion---but did not fully meet all targets and left equity gaps unresolved. In 2025, CSU introduced its first full systemwide strategic plan in more than a decade, paired with a new Student Success Framework and the ``CSU Promise.'' The new agenda keeps graduation important, but expands success to career placement, earnings, socioeconomic mobility, experiential learning, affordability, and alumni-powered career advancement. This evolution is directly relevant to FYE student-employment research because campus jobs are now framed as career-integrated experiential learning that should help students become first-career ready.
\end{abstract}

\sectionrule

\section{Purpose of This Report}
This report was prepared to support the FYE Research \& Engagement project on skill development and career readiness for student employees. It answers three questions:
\begin{enumerate}
  \item What did CSU prioritize over the past five years?
  \item How has strategy changed with \textit{CSU Forward}?
  \item Why does that evolution matter for FYE's weekly competency, feedback, goal-setting, and interview-support practices?
\end{enumerate}

\section{Strategic Context: Why ``Past Five Years'' Means GI~2025}
A key clarification: for most of the last decade, CSU did \textbf{not} operate under a newly refreshed full systemwide strategic plan. The previous major system plan, \textit{Access to Excellence}, dated to approximately 2008. CSU itself describes \textit{CSU Forward} (announced September~9, 2025) as its first systemwide strategic plan in more than a decade.

Therefore, the practical ``past five years'' comparison is:
\begin{itemize}
  \item \textbf{Then:} Graduation Initiative 2025 as the dominant student-success strategy (launched 2016; culminating around 2025).
  \item \textbf{Now:} \textit{CSU Forward} + new Student Success Framework (2025 onward).
\end{itemize}
In plain terms, CSU moved from a powerful \textbf{graduation initiative} to a broader \textbf{system strategy} that includes graduation, career outcomes, affordability, workplace excellence, and statewide impact.

\section{What Came Before: Graduation Initiative Roots}
Before GI~2025, CSU already ran an earlier graduation initiative (beginning around 2009). That earlier effort aimed roughly to:
\begin{itemize}
  \item Raise the six-year first-time graduation rate (from about 46\% toward 54\% by 2015);
  \item Cut racial/ethnic equity gaps in half.
\end{itemize}
Outcomes of that earlier phase were mixed: CSU surpassed the six-year graduation-rate goal, but equity gaps remained nearly unchanged. This set up the next decade's agenda: keep pushing completion, add faster time-to-degree goals, and broaden equity definitions.

\section{Graduation Initiative 2025: The Past-Five-Year Core Agenda}

\subsection{What GI~2025 Tried to Do}
Launched in 2016, GI~2025 was CSU's flagship effort to improve student success through degree completion. Its core logic was:
\begin{quote}
\textit{If more students graduate---and graduate faster---with smaller equity gaps, CSU will better fulfill its public mission.}
\end{quote}

Major systemwide targets included approximately:
\begin{itemize}
  \item 40\% four-year graduation rate for first-time students;
  \item 70\% six-year graduation rate for first-time students;
  \item 45\% two-year graduation rate for transfer students;
  \item 85\% four-year graduation rate for transfer students;
  \item Eliminate equity gaps by race/ethnicity, income (Pell), and first-generation status.
\end{itemize}
Campuses also set campus-specific targets. Funding grew substantially over time (state support plus dedicated tuition revenue), reaching hundreds of millions annually by the final years.

\subsection{How Campuses Worked Toward GI~2025}
CSU guided campus activity through priority areas such as:
\begin{itemize}
  \item Academic preparation (tutoring, readiness supports);
  \item Enrollment management (more course sections; progress pathways);
  \item Student engagement and well-being (including basic needs);
  \item Student financial support;
  \item Data-informed decision making;
  \item Administrative improvements and efficiencies.
\end{itemize}
In practice, campuses used real-time data to identify students off-track and intervene through advising, course access, emergency aid, mentoring, and related supports.

\subsection{What Improved Under GI~2025}
GI~2025 produced clear and important gains:
\begin{itemize}
  \item \textbf{Faster completion:} Four-year first-time graduation rates rose sharply (CSU reports roughly a near-doubling over the initiative window, to about 37\% from about 19\% in commonly cited summaries).
  \item \textbf{Transfer momentum:} Two-year transfer graduation rates improved strongly; CSU met/exceeded the systemwide 45\% target.
  \item \textbf{Scale of impact:} CSU associates the GI~2025 period with roughly 170,000 additional graduates.
  \item \textbf{Campus-level wins:} Multiple campuses posted large gains (e.g., double-digit improvements in four-year and transfer two-year rates).
  \item \textbf{Institutional culture:} Stronger expectations around student success data, transparency, and shared responsibility for completion.
\end{itemize}

\subsection{Where GI~2025 Fell Short}
Independent and CSU-facing assessments also show limits:
\begin{itemize}
  \item Not all system targets were fully met (especially six-year graduation/persistence-related targets).
  \item Equity gaps by race, income, and first-generation status remained wide for many first-time cohorts.
  \item Progress was stronger on helping students finish \textbf{faster} than on graduating substantially more students overall through improved long-term persistence.
  \item Success remained defined mainly as degree completion, with limited system-level tracking of post-graduation career outcomes.
\end{itemize}

\subsection{Bottom Line on the Past Five Years}
The past five years were dominated by a completion agenda. CSU got better at moving students through to degrees---especially more quickly---but still faced unfinished work on equity, long-horizon persistence, and proving career value after graduation.

\section{The Transition: Year of Engagement to a Broader Framework}
As GI~2025 culminated, CSU conducted a broad Year of Engagement involving students, faculty, staff, alumni, and partners. Feedback and research informed a new Student Success Framework designed as the next iteration after GI~2025.

This transition mattered because CSU publicly acknowledged that students' needs had changed: more demand for flexibility, career relevance, holistic support, and clearer pathways into work or further study. The new framework does not discard graduation goals; it builds on them.

\section{CSU Forward (2025): The New Systemwide Strategic Plan}

\subsection{What CSU Forward Is}
\textit{CSU Forward: Thriving Students, Thriving University, Thriving California} is CSU's new system roadmap. It positions CSU as California's socioeconomic mobility engine and seeks to prove the value of public higher education in a tech-driven economy.

\subsection{Four Performance Objectives}
\begin{enumerate}
  \item \textbf{Economic Mobility Engine.} Create flexible, career-integrated learning pathways; strengthen employment and earnings outcomes; expand industry partnerships and adaptive learning.
  \item \textbf{Sustainable Affordability.} Keep high-quality education accessible; pursue debt-free graduation goals for lower-income students; improve operational efficiency and reinvest in students.
  \item \textbf{Employer of Choice.} Strengthen faculty/staff engagement, retention, and professional excellence as drivers of student outcomes.
  \item \textbf{Statewide Impact.} Mobilize alumni and system identity to advance student careers and regional economic development.
\end{enumerate}

\subsection{Example Near-Term Milestones Highlighted by CSU}
\begin{itemize}
  \item Serve about 50,000 learners annually through flexible, modular, career-integrated programs;
  \item Advance debt-free graduation commitments for students from families under \$75,000;
  \item Achieve measurable gains in graduation and employment outcomes;
  \item Improve engagement/retention among CSU's 63,000+ employees;
  \item Mobilize 100,000+ alumni to support student career advancement;
  \item Strengthen CSU's role in regional economic development.
\end{itemize}

\subsection{System KPIs Under CSU Forward}
Progress is framed around indicators such as:
\begin{itemize}
  \item Value of the CSU degree (graduation \textbf{and} employment outcomes);
  \item Affordability and access;
  \item Maximizing California's investment through system scale;
  \item Premier workplace quality;
  \item California's competitive advantage.
\end{itemize}

\section{Student Success Framework and the CSU Promise}
Paired with CSU Forward, the Student Success Framework redefines what ``student success'' means.

\subsection{The CSU Promise}
\begin{quote}
\textit{Every CSU student will graduate with the opportunity for a first career job and/or a clear path to further study.}
\end{quote}

\subsection{Four Commitments}
\begin{enumerate}
  \item Affordability;
  \item Personalized journeys to success;
  \item Experiential learning;
  \item Adaptive and responsive teaching and learning.
\end{enumerate}

\subsection{Expanded Success Metrics}
Beyond retention and graduation, CSU intends to track outcomes such as:
\begin{itemize}
  \item Job or graduate-school placement within 6 and 12 months;
  \item Career earnings at 2, 5, and 10 years;
  \item Experiential learning participation;
  \item Student satisfaction;
  \item Alumni engagement;
  \item Net cost and socioeconomic mobility indicators.
\end{itemize}

\subsection{Why Experiential Learning Matters Here}
The framework treats experiential learning broadly---including internships, undergraduate research, leadership, service learning, and \textbf{on-campus employment}. That is a critical upgrade from older completion-only framing: student jobs are no longer just operational labor; they are potential career-readiness pathways.

\section{Detailed Comparison: Then vs Now}

\begin{center}
\renewcommand{\arraystretch}{1.2}
\begin{tabularx}{\textwidth}{>{\raggedright\arraybackslash}p{0.22\textwidth} >{\raggedright\arraybackslash}X >{\raggedright\arraybackslash}X}
\toprule
\textbf{Dimension} & \textbf{Past $\sim$5 Years (GI~2025)} & \textbf{Now (CSU Forward + Framework)} \\
\midrule
Core question & Did students graduate (on time)? & Did students graduate \textbf{and} become career-ready/mobile? \\
\midrule
Primary metrics & Graduation rates; equity gaps & Graduation + placement + earnings + mobility + satisfaction + experiential learning + alumni engagement \\
\midrule
Student promise & Finish degree; close gaps & First career job and/or further-study path \\
\midrule
Learning emphasis & Courses, advising, basic needs, time-to-degree & Adds career-integrated pathways and experiential learning \\
\midrule
Workforce link & Indirect (finish sooner $\rightarrow$ work sooner) & Direct (skills, networks, employer partnerships, AI career tools) \\
\midrule
Affordability & Important support theme & Central performance objective \\
\midrule
Alumni role & Limited in GI framing & Mobilize 100k+ alumni for career advancement \\
\midrule
Technology & Data for persistence/graduation & Adaptive learning, AI career platforms, tech-ready workforce focus \\
\midrule
Strategy type & Major initiative under shared completion targets & Full system strategic plan with four performance objectives \\
\bottomrule
\end{tabularx}
\end{center}

\section{What Changed and What Improved}

\subsection{Major Changes}
\begin{enumerate}
  \item \textbf{Definition of success expanded.} Completion remains necessary but is no longer sufficient.
  \item \textbf{Career outcomes became explicit.} Placement, earnings, and first-career readiness are now strategic targets.
  \item \textbf{Experiential learning was elevated.} Internships and campus employment are recognized as skill-building systems.
  \item \textbf{Affordability moved to center stage.} Debt and access are treated as structural strategy, not only student-support add-ons.
  \item \textbf{System coordination increased.} CSU emphasizes ``systemness'': shared identity, shared metrics, and cross-campus leverage.
  \item \textbf{Technology strategy matured.} From dashboards for graduation management toward tools that connect skills to employer demand.
\end{enumerate}

\subsection{Clear Improvements in Strategic Design}
\begin{itemize}
  \item Better alignment with public expectations about the value of a college degree;
  \item Stronger link between campus experience and post-graduation mobility;
  \item More honest recognition that persistence/equity work remains unfinished;
  \item Broader accountability (students, employees, alumni, regions);
  \item A more actionable narrative for programs like FYE student employment.
\end{itemize}

\subsection{Continuities (What Did Not Disappear)}
CSU did not abandon GI~2025 priorities. Graduation rates, retention, equity-minded support, advising, and basic-needs concerns remain relevant. New graduation/retention targets are still being set at ambitious levels (for example, higher four- and six-year goals discussed in transition materials). The shift is expansion, not replacement.

\section{Implications for the FYE Project}
The FYE project---assessing weekly competency reflection, assessment/feedback, goal-setting, and interview support for student employees---maps directly onto the strategic evolution described above.

\subsection{Why FYE Work Matters More Under CSU Forward}
\begin{itemize}
  \item FYE student employment is experiential learning, which the new framework explicitly values;
  \item Weekly Practice Plans, Self-Reports, feedback, and 1:1/2:1 structures support personalized journeys;
  \item Interview readiness and evidence-based skill stories support the CSU Promise of first-career opportunity;
  \item Light AI/tech readiness aligns with CSU's tech-driven workforce focus;
  \item Alumni/post-graduate cohort inquiry aligns with alumni-powered career advancement goals.
\end{itemize}

\subsection{Recommended Framing Language for Supervisors/Admins}
\begin{quote}
\textit{Over the past five years, CSU strategy centered on Graduation Initiative 2025 and improved degree completion. CSU Forward evolves that agenda by adding career readiness, experiential learning, and post-graduation mobility. This FYE project supports that evolution by assessing whether student-employment practices build career-ready skills and by identifying improvements in interview readiness, evidence-based skill development, and tech/AI literacy.}
\end{quote}

\section{Conclusions}
\begin{enumerate}
  \item The past five years were defined primarily by GI~2025's completion agenda.
  \item GI~2025 delivered real improvement in faster graduation and transfer completion, with incomplete equity and persistence outcomes.
  \item \textit{CSU Forward} is a strategic upgrade: from ``graduate more students'' to ``graduate career-ready students who can thrive.''
  \item The Student Success Framework operationalizes that upgrade through the CSU Promise and expanded metrics.
  \item FYE's weekly developmental employment model is highly aligned with the new agenda and is therefore a strong site for assessment and improvement.
\end{enumerate}

\section{Sources Consulted}
\begin{itemize}
  \item CSU news: \textit{CSU Sets the Standard... New Strategic Plan} (Sept.~9, 2025).\\
  \url{https://www.calstate.edu/csu-system/news/Pages/CSU-systemwide-strategic-plan-2025.aspx}
  \item CSU Strategic Planning hub: \textit{CSU Forward}.\\
  \url{https://www.calstate.edu/Strategic-Planning-for-the-CSU}
  \item CSU Forward plan PDF (system strategic plan materials, 2025).
  \item CSU Student Success Framework final report / Year of Engagement materials (2025).
  \item CSU news: \textit{CSU Culminates 10-Year Graduation Initiative...} (Oct.~17, 2025).
  \item Legislative Analyst's Office: \textit{Assessment of CSU's Graduation Initiative 2025} (Dec.~10, 2025).
\end{itemize}

\vspace{1.2em}
{\footnotesize\textit{Note: Figures cited from CSU and LAO public summaries are rounded for readability. For formal publication or grant use, verify the latest dashboard values and cite exact cohort years.}}

\end{document}

